\newpage
\null
\thispagestyle{empty}
\newpage
%\begin{refsection}%[chapters/Ref/main123fixed.bib]

\vspace*{\fill}
\begin{center}
	\chapterpart{\Large\texorpdfstring{\centering Paper 1\\[1.5em]
			{\large Published in \textit{Journal}, Paper}\\[1em]
			Paper 1 reference}{Paper 1}}
\end{center}
\vspace*{\fill}
\clearpage

\section{Abstract}
Angiogenesis is a process for the formation of new blood vessels from existing vessels and it sits directly at the intersection of normal physiology and multiple diseases.
Normal physiological activities like embryonic development, tissue homoeostasis require precise regulation.


\subsection{Statistical Analysis}	
Detailed ANOVA and Statistical analysis method is mentioned in section~\ref{Statistical} on page~\pageref{Statistical}.
\section{Results and Discussion}

This how to cite figure. %(Fig.~\ref{fig:Angio1})
Normalization of images redistributes pixel intensities to a consistent dynamic range and helps in suppression of background gradients while vessel contrast is preserved.
ROI cropping helps in elimination of peripheral tissue and non-vascular background 
\clearpage
%\printbibliography
\clearpage
%\end{refsection}